\documentclass[twocolumn,a4paper]{article}
\usepackage{graphicx} % Required for inserting images
\usepackage{hyperref} % Required for \url command

\title{Simulation of a computer game by diffusion model}
\author{
  Lukáš Foukal \and
  David Nepraš \and
  Michael Peštuka
}
\date{March 2026}

\begin{document}

\maketitle

\textit{On 13.3. you have to submit 1 A4 in the system. Describe the problem you are solving, the approaches you draw inspiration from, data sets you are planning to use and a plan for your solution and experiments. You also have to include a link to a public GIT repository containing your project source codes. During the following week (16.-20.2.) you have to schedule a consultation with your project mentor to discuss the proposal.}

\section{Problem}
The goal of our project is to create a playable simulation of a computer game by using a diffusion model. For such a model, the input is the previous frame (or a sequence of frames) and the user action. The output is the next frame.

\section{Data}
It is our intention to generate our own data using pre-existing games.


\section{Approach and Experiments}

We have drawn inspiration from \textit{GameNGen}\footnote{https://arxiv.org/pdf/2408.14837} and \textit{Oasis}\footnote{https://oasis-model.github.io/}. Both comprise of an auto-encoder and a diffusion model. The former is based on pre-trained \textit{Stable Diffusion 1.4} models while the latter uses a transformer-based architecture with \textit{ViT}\footnote{https://arxiv.org/pdf/2010.11929} and \textit{DiT}\footnote{https://arxiv.org/pdf/2212.09748}.

In Appendix A of \textit{GameNGen}, section 10, the authors describe their inital experiment on the dinosaur game present in Google Chrome web browser, noting that their approach worked well with only 2000 episodes of training data and 3000 training epochs.This shows us the feasibility of scaling down their approach. The fact that pre-trained models are used and only fine tuned is also beneficial.

\section{GIT Repository}

\url{https://github.com/MichaelPestuka/KNN_proj}

\end{document}

